\section{Introduction}
3D scene understanding supports question answering, object description, and language grounding in spatial observations. Models such as LLaVA-3D, 3D-LLaVA, and LEO~\cite{zhu2024llava3d,deng20253dllava,huang2024leo} connect multi-view or point-cloud representations to a language model (LM). These 3D large multimodal models (3D LMMs) process hundreds to thousands of spatial tokens, increasing prefill computation and decoding key--value (KV) storage. We study token reduction for efficient inference without retraining these models.

Existing methods address two coupled choices: which tokens to retain and when to remove the rest. FastV uses last-prompt attention to rank visual tokens~\cite{chen2024fastv}. Depth-aware approaches also organize removal across layers: PyramidDrop exploits depth-dependent redundancy, while HiMAP uses information flow to guide hierarchical pruning~\cite{yuan2024pyramiddrop,yin2025himap}. Alongside these scheduling strategies, 3DZip uses spatial structure and feature diversity to compress 3D inputs~\cite{baek2026zip}. Together, these approaches motivate studying token selection and pruning depth jointly.

For 3D LMMs, the central question is whether delaying token removal improves the quality--cost trade-off under a fixed scoring rule. Applying the same rule at different layers can still select different tokens, coupling pruning depth with token identity. Evaluating a schedule therefore requires testing both where tokens are removed and whether the retained set supports task performance. Analytical cost estimates must also be distinguished from application time, since equal retained-token counts need not imply equal computational or runtime savings.

Across two Llama-family 7B models with image-derived and superpoint representations, we observe low intermediate spatial attention followed by a concentrated late resurgence. This observation motivates \textbf{GeoGate}, a training-free schedule for late-layer token pruning. GeoGate retains the full sequence through zero-indexed block L14 and applies FastV's existing ranking before L15. The observed attention pattern guides the schedule; controlled experiments assess its effect on task quality.

GeoGate reuses the pruning block's full-length forward pass to obtain rankings, avoiding an additional scoring pass. It preserves all text tokens and original rotary positions, then compacts every layer's prompt KV cache to the retained set before decoding. This synchronization reduces prompt storage and supports the tested decode interface. L14 and the default 10\% retention budget were selected on ScanQA validation. Retention tests assess retained-set quality, and fixed-set deletion controls hold token identity constant across depths. Task transfer and application timing assess the quality and cost of this combined prefill/cache intervention.

Our contributions are:
\begin{enumerate}
\item \textbf{A training-free late-layer pruning schedule.} GeoGate provides a one-pass prefill schedule for late-layer spatial-token pruning with FastV's existing scorer. On LLaVA-3D ScanQA validation, L14 with 128 retained spatial tokens exceeds selected L2/L8 controls by 1.22/1.03 exact-match (EM) points at matched analytical LM cost.
\item \textbf{Controlled evidence on deletion depth.} Deletion depth affects task performance even when the deleted token set is held fixed. For sets selected at L14 and L20, deletion causes greater loss relative to the unpruned model at L8 than at L14. This ordering supports a depth effect without identifying L14 as the most damaging deletion location.
\item \textbf{Task quality and practical efficiency.} Task-transfer tests reveal quality losses despite favorable ScanQA point estimates. For LLaVA-3D at 10\% retention, analytical LM FLOPs decrease by 42.1\%, while application time decreases by 0.8--1.4\% on one H100.
\end{enumerate}
